\chapter{Gauge Theory Of Deforming Bodies}
\label{chap:deforming-body}

Rigid-body mechanics separates orientation from shape because a rigid body has
no internal shape variables.  A deforming body has both.  Frank Wilczek's
gauge-theoretic viewpoint, developed with Shapere in the context of swimming
and self-propulsion,\footnote{A. Shapere and F. Wilczek, ``Gauge kinematics of
deformable bodies,'' \emph{American Journal of Physics} 57, 514--518 (1989),
\href{https://doi.org/10.1119/1.15986}{doi:10.1119/1.15986}. Their closely
related low-Reynolds-number swimming formulation appears in
\href{https://doi.org/10.1017/S002211208900025X}{doi:10.1017/S002211208900025X}.}
treats the decomposition into overall rigid motion and internal deformation as
a choice of gauge.  The physically meaningful object is not the instantaneous
frame itself, but the connection that tells us which rigid motion accompanies an
infinitesimal change of shape under a specified momentum constraint.

The central phenomenon is geometric phase.  A body can execute a closed loop in
shape space and return to the same internal shape while acquiring a net rotation
or translation in physical space.  This is the mechanical analogue of parallel
transport in a gauge field.  The curvature of the connection measures the
leading displacement produced by small loops, while the full holonomy records
the finite motion produced by large cycles.

\paragraph{Roadmap.}
The chapter starts with the finite-dimensional Shapere-Wilczek split between
rigid placement \(g\) and shape \(s\). It then derives the mechanical
connection from zero momentum and explains how this connection transforms under
changes of body frame. The two-rotor model is a minimal laboratory for
noncommuting connection components. The last third returns to continuum
mechanics: local frames, Cosserat rods, curvature, torsion, and defects show
how the same gauge language becomes a field theory over a material body.

\paragraph{Reading rule.}
There are two meanings of ``gauge'' in this chapter. One is finite-dimensional:
choosing body axes over shape space. The other is local and continuum: choosing
a frame at each material point. Both are frame choices, and both have
curvatures, but they live over different base spaces. Whenever a symbol
\(\mathcal A\), \(A\), \(\omega\), or \(R\) appears, first identify whether the
base variable is shape \(s\), time \(t\), or material position \(X\).

\section{Gauge Kinematics Of Deformable Bodies}
\label{sec:gauge-kinematics-deformable-bodies}

The preceding chapter described how an elastic body deforms relative to a
reference body. Gauge kinematics asks a different but compatible question: if
the body changes shape internally, how much of its observed motion is a genuine
rigid displacement, and how much is a convention for choosing body axes at each
shape? This is the same distinction that appears in gauge theory: local
descriptions depend on a frame choice, while holonomy around a closed loop is
observable.

The Wilczek-Shapere point of view begins with a finite deformable body moving in
space. Separate its configuration into:
\[
  g\in G,\qquad s\in\mathcal S,
\]
where \(G\) is the rigid-motion group, usually \(\SE(3)\) or \(\SO(3)\), and
\(\mathcal S\) is shape space. A point of the full configuration space is a
rigid placement \(g\) together with an internal shape \(s\).

The subtlety is that the split into \(g\) and \(s\) is not unique. To say what
the body's orientation is at a given shape, one must choose body axes for that
shape. A shape-dependent change of axes
\[
  h:\mathcal S\to G
\]
changes the representative \(g\). This is exactly a gauge choice.

Geometrically, the full configuration space is locally a product
\[
  Q_{\mathrm{full}}\simeq G\times\mathcal S,
\]
but no preferred product splitting is supplied by physics. The projection
\[
  Q_{\mathrm{full}}\to\mathcal S
\]
forgets the rigid placement and remembers only shape. A connection is a rule
for deciding which infinitesimal motions in \(Q_{\mathrm{full}}\) count as
``pure shape change'' and which contain rigid motion. In mechanics this rule is
not arbitrary: the mechanical connection is selected by the kinetic energy and
the momentum constraint.

The observable content is therefore not the connection coefficient by itself.
A different convention for body axes changes the coefficient, just as a
different electromagnetic gauge changes the vector potential. What survives is
the reconstructed rigid motion, especially the holonomy around a closed shape
cycle and the curvature that controls small cycles. This is the sense in which
the Shapere-Wilczek description is a gauge theory of deforming bodies.

\subsection{Principal-Bundle Form Of The Split}
\label{subsec:deforming-body-principal-bundle}

The clean mathematical model is a principal bundle.  Let \(Q_{\mathrm{full}}\)
be the full configuration space of the deforming body, after any constraints
such as fixed center of mass have been imposed.  The rigid-motion group \(G\)
acts on \(Q_{\mathrm{full}}\) by changing the overall placement of the same
internal shape.  The shape space is the quotient
\[
  \mathcal S=Q_{\mathrm{full}}/G,
  \qquad
  \pi:Q_{\mathrm{full}}\to\mathcal S .
\]
The point \(s=\pi(q)\) remembers the shape and forgets the rigid placement.

\begin{assumption}[Regular shape bundle]
In this chapter the \(G\)-action is assumed to be free and proper on the
configurations under discussion, so that \(\mathcal S\) is a smooth manifold
and \(Q_{\mathrm{full}}\to\mathcal S\) is locally a product.  Shapes with
extra symmetry, such as a perfectly spherical body, have stabilizers and require
minor orbifold or stratified-space modifications.  The formulas below describe
the generic regular stratum.
\end{assumption}

A local choice of body frame is a local section
\[
  \sigma:\mathcal U\subset\mathcal S\to Q_{\mathrm{full}}.
\]
Every nearby configuration can then be written as
\[
  q=g\,\sigma(s),
  \qquad g\in G,\quad s\in\mathcal U.
\]
This expression is not unique because the section is not unique.  If the
internal representative is changed by a shape-dependent rigid motion
\[
  \sigma'(s)=h(s)^{-1}\sigma(s),
  \qquad h:\mathcal U\to G,
\]
then the same physical configuration is represented by
\[
  q=g\,\sigma(s)=g'\sigma'(s),
  \qquad
  g'=g\,h(s).
\]
Thus a change of section changes the coordinate \(g\) even though the physical
configuration \(q\) has not changed.  This is the precise analogue of a gauge
transformation.

A connection is a \(G\)-equivariant rule that splits a tangent vector
\(\dot q\in T_qQ_{\mathrm{full}}\) into a vertical rigid part and a horizontal
shape-changing part.  In the local coordinates \((g,s)\), write the body-frame
group velocity as
\[
  \xi=g^{-1}\dot g\in\mathfrak g .
\]
The horizontal condition has the form
\[
  \xi+\mathcal A_\alpha(s)\dot s^\alpha=0.
\]
Therefore a horizontal lift of a shape path obeys the reconstruction equation
\[
  g^{-1}\dot g=-\mathcal A_\alpha(s)\dot s^\alpha .
\]
The local coefficient \(\mathcal A=\mathcal A_\alpha ds^\alpha\) is the gauge
potential.  Under the change of section above it transforms as
\[
  \mathcal A'
  =
  \Ad_{h^{-1}}\mathcal A-h^{-1}\dd h,
\]
which is exactly the inhomogeneous transformation law expected of a connection.
The curvature and holonomy are developed below; the important point here is
that the connection is the mathematical object that makes the phrase ``pure
shape change'' meaningful.

\begin{figure}[htbp]
\centering
\begin{tikzpicture}[scale=1.0, >=Latex, font=\small]
  \draw[->] (-0.2,0) -- (6.5,0) node[right] {\(s^1\)};
  \draw[->] (0,-0.2) -- (0,4.1) node[above] {\(s^2\)};
  \draw[thick, blue, ->] (1.2,1.2) .. controls (2.8,3.4) and (5.0,3.1) .. (5.0,1.4);
  \draw[thick, blue, ->] (5.0,1.4) .. controls (3.9,0.4) and (2.2,0.3) .. (1.2,1.2);
  \node[blue, fill=white, inner sep=2pt] at (3.25,3.58) {shape cycle};
  \begin{scope}[shift={(1.2,1.2)}, scale=0.28]
    \draw[fill=gray!20, thick] (-0.8,0) ellipse (0.45 and 0.25);
    \draw[fill=gray!20, thick] (0.8,0) ellipse (0.45 and 0.25);
    \draw[thick] (-0.35,0) -- (0.35,0);
  \end{scope}
  \begin{scope}[shift={(5.0,1.4)}, scale=0.28, rotate=-20]
    \draw[fill=gray!20, thick] (-1.0,0) ellipse (0.55 and 0.18);
    \draw[fill=gray!20, thick] (1.0,0) ellipse (0.55 and 0.18);
    \draw[thick] (-0.45,0) -- (0.45,0);
  \end{scope}
  \begin{scope}[shift={(5.35,3.05)}, scale=0.28, rotate=45]
    \draw[fill=gray!20, thick] (-0.8,0) ellipse (0.45 and 0.25);
    \draw[fill=gray!20, thick] (0.8,0) ellipse (0.45 and 0.25);
    \draw[thick] (-0.35,0) -- (0.35,0);
    \draw[red!80!black, thick, ->] (0.0,0.65) arc (90:300:0.45);
  \end{scope}
  \draw[red!80!black, thick, ->] (4.85,2.72) -- (5.18,2.92);
  \node[align=center, fill=white, inner sep=2pt] at (2.55,0.62)
    {closed shape loop\\can have nonzero holonomy};
  \node[red!80!black, align=center, fill=white, inner sep=2pt] at (5.75,2.35)
    {net rigid\\rotation};
\end{tikzpicture}
\caption{Gauge kinematics: a closed loop in shape space can produce a net rigid
motion. The start and end shape can be the same while the rigid placement has
changed; the effect is the holonomy of a connection.}
\label{fig:gauge-shape-cycle-holonomy}
\end{figure}

\subsection{Mechanical Connection From Zero Momentum}
\label{subsec:mechanical-connection-zero-momentum}

The mechanical connection is obtained by solving a constraint, not by adding a
new force law. If the body has no total rigid momentum, then any internal shape
velocity that would otherwise create angular or linear momentum must be
cancelled by an accompanying rigid velocity. The connection is the linear rule
that performs this cancellation at each shape.

The zero-momentum assumption is a modelling choice, not a mathematical
necessity. It is appropriate for problems such as a freely falling cat or an
isolated rotor device initialized with no angular momentum. With nonzero
momentum, the same kinetic metric still defines a connection, but the
reconstruction equation contains an additional momentum-driven drift. The
zero-momentum case is the cleanest place to see geometric phase because the
entire rigid motion is produced by shape change.

Let the body velocity in the moving frame be
\[
  \xi=g^{-1}\dot g\in\mathfrak g,
\]
where \(\mathfrak g\) is the Lie algebra of \(G\). Suppose the kinetic energy
has the block form
\[
  T
  =
  \frac12
  \begin{pmatrix}
  \xi\\ \dot s
  \end{pmatrix}^T
  \begin{pmatrix}
  M(s) & A(s)\\
  A(s)^T & m(s)
  \end{pmatrix}
  \begin{pmatrix}
  \xi\\ \dot s
  \end{pmatrix}.
\]
This is not an additional assumption about linearity of the equations of
motion; it is simply the quadratic form of kinetic energy written in the
variables \((\xi,\dot s)\).  The block \(M(s)\) is the locked inertia operator:
it measures the rigid momentum produced by a rigid velocity when the shape is
held fixed.  The block \(m(s)\) is the shape-space metric for pure deformation
in the chosen frame.  The off-diagonal block \(A(s)\) is the crucial one: it
measures the rigid momentum generated by changing shape while the chosen body
frame is held fixed.

\begin{assumption}[Regular locked inertia]
The zero-momentum reconstruction formula below assumes that \(M(s)\) is
invertible on the rigid motions being reconstructed.  If a shape has an exact
symmetry that makes some rigid velocity physically invisible, then \(M(s)\)
has a null direction and the regular principal-bundle description must be
restricted or quotient further.
\end{assumption}

The momentum conjugate to rigid motion is
\[
  \Pi=\frac{\partial T}{\partial \xi}=M(s)\xi+A(s)\dot s.
\]
The symbol \(\Pi\) is the total rigid-motion momentum written in body
variables: for \(G=\SE(3)\) it packages total linear and angular momentum.
If the total angular and linear momentum vanish,
\[
  \Pi=0,
\]
as they do for an isolated body prepared with zero net external impulse, then
\(\Pi\) remains zero by conservation and
\[
  \xi=-M(s)^{-1}A(s)\dot s.
\]
Writing \(\dot s=\dot s^\alpha\partial_\alpha\), define the connection
\[
  \boxed{
  \mathcal A_\alpha(s)=M(s)^{-1}A_\alpha(s).
  }
\]
The rigid velocity is then
\[
  \boxed{
  g^{-1}\dot g=-\mathcal A_\alpha(s)\dot s^\alpha.
  }
\]
This is the cleanest derivation of the gauge field: it is the rule that tells a
zero-momentum deforming body how much rigid motion is forced by an infinitesimal
shape change.

The minus sign is part of the reconstruction convention. A positive connection
component means that the chosen shape change would create rigid momentum in the
body frame; the actual zero-momentum motion supplies the opposite rigid
velocity to cancel it. This is the mechanical meaning behind the gauge-theory
formula.

\subsection{Explicit Inertia-Operator Form}
\label{subsec:explicit-gauge-equation}

A concrete \(\SO(3)\) version of the same construction is useful because it
turns the abstract mechanical connection into a linear algebra problem. Let
\[
  A=R^{-1}\dot R\in\mathfrak{so}(3)
\]
be the angular velocity matrix of the body frame. Let \(Q=Q^T\) encode the
instantaneous shape through the second moment tensor
\[
  Q_{ij}=\int_B\rho(X)X_iX_j\,\dd^3X
\]
in the chosen internal frame. The corresponding inertia tensor is
\[
  \widetilde I_{ij}
  =
  \delta_{ij}\sum_{k=1}^3Q_{kk}-Q_{ij}.
\]
Let \(Y(X,s(t))\) be the position of a material point in the chosen internal
frame, with the center of mass at the origin. The spatial position is
\(x=RY\). In the body frame,
\[
  R^{-1}\dot x=A Y+\dot Y.
\]
The first term is the rigid velocity required by the frame motion; the second
term is internal deformation. The body angular momentum is therefore
\[
  M_{\mathrm{tot}}
  =
  \int_B\rho\,Y\times(AY+\dot Y)\,\dd^3X.
\]
With the hat convention used for the rigid body,
\(\widehat v\,w=v\times w\), or
\[
  (\widehat v)_{ij}=-\epsilon_{ijk}v_k,
\]
the rigid part obeys the matrix identity
\[
  \widehat{\widetilde I\Omega}=AQ+QA,
  \qquad A=\widehat\Omega.
\]
The internal part is
\[
  J_{\mathrm{int}}
  =
  \int_B\rho\,Y\times\dot Y\,\dd^3X.
\]
For zero total angular momentum,
\[
  \widetilde I\Omega+J_{\mathrm{int}}=0.
\]
Define the deformation-generated momentum that the rigid motion must cancel by
\[
  \widetilde J=-J_{\mathrm{int}}.
\]
Then the zero-momentum condition takes the linear matrix form
\[
  \boxed{
  AQ+QA=\widehat{\widetilde J}.
  }
\]
Solving this equation is the explicit finite-dimensional gauge connection:
the angular velocity \(A\) is determined by the shape velocity through
\(\widetilde J\).

\begin{derivation}[Solving the finite-dimensional gauge equation]
Write
\[
  A=\widehat\Omega,\qquad A_{ij}=-\epsilon_{ijk}\Omega_k.
\]
To verify the identity \(AQ+QA=\widehat{\widetilde I\Omega}\), diagonalize the
symmetric tensor \(Q\) at the instant under consideration:
\[
  Q=\diag(q_1,q_2,q_3).
\]
Then
\[
  (AQ+QA)_{12}=-(q_1+q_2)\Omega_3,\quad
  (AQ+QA)_{23}=-(q_2+q_3)\Omega_1,\quad
  (AQ+QA)_{31}=-(q_3+q_1)\Omega_2.
\]
But the inertia tensor has principal moments
\[
  \widetilde I_1=q_2+q_3,\qquad
  \widetilde I_2=q_3+q_1,\qquad
  \widetilde I_3=q_1+q_2,
\]
so these entries are exactly those of
\(\widehat{\widetilde I\Omega}\). Since both sides transform tensorially under
orthogonal changes of the internal frame, the identity holds in any frame.
At zero total angular momentum, the rigid angular momentum must equal
\(\widetilde J\), the negative of the internal contribution. The linear map
from angular velocity \(\Omega\) to rigid angular momentum is precisely the
inertia tensor:
\[
  \widetilde J_i=\widetilde I_{ij}\Omega_j.
\]
Equivalently,
\[
  \Omega_i=(\widetilde I^{-1})_{ij}\widetilde J_j.
\]
Therefore
\[
  \boxed{
  A_{ij}
  =
  -\epsilon_{ijk}
  (\widetilde I^{-1})_{k\ell}\widetilde J_\ell.
  }
\]
This is the concrete \(\SO(3)\) deforming-body gauge formula. In the abstract
notation above, \(M(s)\) is the inertia operator
\(\widetilde I(s)\), while the cross term \(A(s)\dot s\) is the internally
generated angular momentum \(\widetilde J(s,\dot s)\). The helicopter rotor
laboratory in Section~\ref{sec:lab-helicopter-gauge} is the finite-dimensional
special case where the shape variables are two rotor angles and this abstract
inertia-operator equation reduces to a pair of scalar connection components.
\end{derivation}

\subsection{Gauge Transformation And Holonomy}
\label{subsec:gauge-holonomy}

Under a shape-dependent change of body frame \(h(s)\), the connection transforms
as
\[
  \boxed{
  \mathcal A'
  =
  \Ad_{h^{-1}}\mathcal A-h^{-1}\dd h.
  }
\]
The sign of the inhomogeneous term follows from the reconstruction convention
\(g^{-1}\dot g=-\mathcal A(\dot s)\). With the opposite convention, the same
geometric connection is written with the more familiar plus sign. The base space
is shape space and the structure group is the group of rigid motions.

For a path \(s(t)\), the rigid motion is obtained by integrating
\[
  g^{-1}\dot g=-\mathcal A_\alpha(s)\dot s^\alpha.
\]
For a closed shape cycle \(\Gamma\), the net displacement is
\[
  g(T)
  =
  g(0)\,
  \mathcal P\exp\left(
    -\oint_\Gamma \mathcal A
  \right),
\]
where \(\mathcal P\) denotes path ordering. The result can be nontrivial even
though the shape returns to its starting point.

The curvature of the connection is
\[
  \boxed{
  \mathcal F
  =
  \dd\mathcal A+\mathcal A\wedge\mathcal A,
  }
\]
or in components,
\[
  \mathcal F_{\alpha\beta}
  =
  \partial_\alpha\mathcal A_\beta
  -
  \partial_\beta\mathcal A_\alpha
  +
  [\mathcal A_\alpha,\mathcal A_\beta].
\]

\begin{derivation}[Small-loop holonomy]
Take shape coordinates \((s^1,s^2)\) and traverse a tiny rectangle based at
\(s_0\): first \(+\Delta s^1\), then \(+\Delta s^2\), then
\(-\Delta s^1\), then \(-\Delta s^2\). Along each side the reconstruction
equation
\[
  g^{-1}\dot g=-\mathcal A_\alpha\dot s^\alpha
\]
produces an infinitesimal group exponential. If \(\mathcal A\) were constant,
the product would be the group commutator
\[
  e^{-\mathcal A_2\Delta s^2}
  e^{-\mathcal A_1\Delta s^1}
  e^{\mathcal A_2\Delta s^2}
  e^{\mathcal A_1\Delta s^1}
  =
  \exp\left(
    -[\mathcal A_1,\mathcal A_2]\Delta s^1\Delta s^2
    +O(|\Delta s|^3)
  \right),
\]
with the sign following the reconstruction convention. If \(\mathcal A\) varies
over the rectangle, the two \(s^1\)-edges differ by
\(\partial_2\mathcal A_1\Delta s^2\), and the two \(s^2\)-edges differ by
\(\partial_1\mathcal A_2\Delta s^1\). Combining the derivative terms with the
commutator gives
\[
  g_{\mathrm{loop}}
  =
  \exp\left(
    -\mathcal F_{12}\Delta s^1\Delta s^2
    +O(|\Delta s|^3)
  \right).
\]
This is the finite-dimensional version of the usual curvature-as-holonomy
calculation for a gauge connection.
\end{derivation}

For a small loop enclosing area element
\(\Delta S^{\alpha\beta}\),
\[
  \Delta g
  \approx
  \exp\left(
    -\mathcal F_{\alpha\beta}\Delta S^{\alpha\beta}
  \right).
\]

\begin{figure}[htbp]
\centering
\begin{tikzpicture}[scale=1.0, >=Latex, font=\small]
  \draw[->] (-0.2,0) -- (6.9,0) node[right] {\(s^1\)};
  \draw[->] (0,-0.2) -- (0,3.75) node[above] {\(s^2\)};
  \coordinate (a) at (1.35,1.25);
  \coordinate (b) at (3.65,1.25);
  \coordinate (c) at (3.65,2.55);
  \coordinate (d) at (1.35,2.55);
  \fill[blue!10] (a) -- (b) -- (c) -- (d) -- cycle;
  \draw[blue!70!black, thick, ->] (a) -- node[below] {\(\delta s^1\)} (b);
  \draw[blue!70!black, thick, ->] (b) -- node[right] {\(\delta s^2\)} (c);
  \draw[blue!70!black, thick, ->] (c) -- (d);
  \draw[blue!70!black, thick, ->] (d) -- (a);
  \draw[red!75!black, thick, ->] (2.35,2.15) arc (210:-80:0.42);
  \draw[red!75!black, ->] (4.45,2.55) -- (2.95,2.22);
  \node[red!75!black, align=left, fill=white, inner sep=2pt] at (5.65,2.78)
    {net rigid motion\\\(-\mathcal F_{12}\delta s^1\delta s^2\)};
  \node[align=center, fill=white, inner sep=2pt] at (2.5,0.50)
    {same infinitesimal shape loop,\\different final placement};
\end{tikzpicture}
\caption{Curvature is the local density of holonomy. Traversing a small
rectangle in shape space returns to the same shape but can leave a residual
rigid displacement proportional to the curvature through the enclosed area.}
\label{fig:gauge-curvature-small-loop}
\end{figure}

Thus curvature is the local density of geometric phase. This is the same
mathematical idea as Hannay angle, Berry phase, and the rigid-body falling-cat
effect, now written for deforming bodies.

Flat connections can still produce frame changes along open paths, but a small
closed loop detects curvature. This is why the curvature, rather than the
connection coefficients themselves, is the local observable of the gauge
description. Different choices of body frame change \(\mathcal A\), while the
holonomy around a loop changes by conjugation and therefore represents the same
physical displacement.

\section{A Two-Rotor Deforming-Body Model}
\label{sec:two-rotor-deforming-body-model}

A minimal finite-dimensional example is a rigid carrier with two internal
rotors. The rotor angles are shape variables
\[
  s=(\alpha,\beta)\in S^1\times S^1.
\]
This is the compact version of the helicopter-rotor laboratory in
Section~\ref{sec:lab-helicopter-gauge}: here we strip the device down to its
two noncommuting shape directions so the curvature calculation can be seen in
one page.
Let the carrier have isotropic inertia \(I_0\mathbf 1\). Let one rotor spin
about the body \(z\)-axis with inertia \(I_1\), and another spin about the
body \(y\)-axis with inertia \(I_2\). If \(\Omega\) is the carrier's body
angular velocity, the rotational kinetic energy is
\[
  T=
  \frac12I_0\norm{\Omega}^2
  +
  \frac12I_1(\Omega_z+\dot\alpha)^2
  +
  \frac12I_2(\Omega_y+\dot\beta)^2.
\]
The body angular momentum conjugate to \(\Omega\) is
\[
  \Pi
  =
  I_0\Omega
  +
  I_1(\Omega_z+\dot\alpha)e_z
  +
  I_2(\Omega_y+\dot\beta)e_y.
\]
Imposing zero total angular momentum gives
\[
  \Omega
  =
  -a\,\dot\alpha\,e_z
  -
  b\,\dot\beta\,e_y,
  \qquad
  a=\frac{I_1}{I_0+I_1},
  \qquad
  b=\frac{I_2}{I_0+I_2}.
\]
Thus
\[
  \mathcal A_\alpha=a\,e_z,
  \qquad
  \mathcal A_\beta=b\,e_y.
\]
The connection coefficients are constant, so the derivative part of curvature
vanishes. But \(\SO(3)\) is nonabelian:
\[
  [e_z,e_y]=e_z\times e_y=-e_x.
\]
Therefore
\[
  \mathcal F_{\alpha\beta}
  =
  [\mathcal A_\alpha,\mathcal A_\beta]
  =
  -ab\,e_x.
\]
A small rectangular stroke with sides \(\Delta\alpha,\Delta\beta\) produces a
body rotation whose leading rotation vector is proportional to
\[
  ab\,\Delta\alpha\,\Delta\beta\,e_x
\]
up to the sign convention set by the direction of traversal and by
\(g^{-1}\dot g=-\mathcal A(\dot s)\). The point is not the sign. The point is
that a closed loop in internal shape variables produces a rigid rotation even
though the rotors return to their starting angles. This is the same
geometric-phase mechanism as the falling cat, written as a connection over
shape space.

\section{Swimming At Low Reynolds Number}
\label{sec:low-reynolds-swimming}

Shapere and Wilczek's low-Reynolds-number swimming problem gives a particularly
physical gauge theory. At very small Reynolds number, inertia drops out of
Navier-Stokes and the surrounding fluid satisfies Stokes equations:
\[
  -\nabla p+\mu\Delta u=0,\qquad \nabla\cdot u=0.
\]
Boundary velocity is prescribed by the swimmer's rigid motion plus deformation
rate. Because the Stokes problem is linear, the force and torque on the body
depend linearly on:
\[
  \xi=g^{-1}\dot g
  \qquad\text{and}\qquad
  \dot s.
\]
A freely swimming body has zero net force and torque, so the same algebra as
above gives
\[
  \xi=-\mathcal A_\alpha(s)\dot s^\alpha.
\]
The connection is now computed from a hydrodynamic boundary-value problem
rather than from an inertia tensor. A reciprocal one-parameter stroke encloses
no area in shape space and produces no net motion: this is Purcell's scallop
theorem in gauge language. A two-parameter stroke can enclose area and swim by
curvature.

\section{Local Frames In Elastic Continua}
\label{sec:local-frames-elastic-continua}

The finite-dimensional gauge picture becomes a field theory when every material
point carries a local frame. A Cosserat body has, in addition to
\(\varphi(X,t)\), a rotation field
\[
  R(X,t)\in\SO(3).
\]
The directors are
\[
  d_a(X,t)=R(X,t)E_a,
\]
where \(E_a\) is a fixed reference basis. The rotational connection is
\[
  A_A=R^{-1}\partial_A R\in\mathfrak{so}(3),
  \qquad
  A_t=R^{-1}\partial_t R.
\]
The translational strain measured in the local frame is
\[
  \Gamma_A=R^{-1}\partial_A\varphi.
\]
For a rigid body, \(R\) depends only on time. For a Cosserat continuum, \(R\)
depends on \(X\), so neighboring material points can rotate relative to one
another.

This is the field-theoretic version of the rigid-body frame issue. In the
rigid body there is one moving frame; in a Cosserat body there is a moving frame
at every material point. The connection \(A_A\) compares nearby frames, while
\(A_t\) records how a given local frame evolves in time.

\begin{definition}[Local frame gauge transformation]
Let \(G(X,t)\in\SO(3)\) be a change of local material frame:
\[
  R'(X,t)=R(X,t)G(X,t).
\]
Then
\[
  A'_A=G^{-1}A_AG+G^{-1}\partial_A G,
  \qquad
  \Gamma'_A=G^{-1}\Gamma_A.
\]
Thus \(A_A\) is a gauge connection and \(\Gamma_A\) is a vector-valued
matter field in the local frame.
\end{definition}

The simplest elastic energy for a Cosserat continuum penalizes deviations of
these gauge-covariant strains from preferred values:
\[
  W_{\mathrm{Cosserat}}
  =
  \frac12 C^{AB}_{ab}
  \left(\Gamma_A^a-\Gamma_{0A}^a\right)
  \left(\Gamma_B^b-\Gamma_{0B}^b\right)
  +
  \frac12 D^{AB}_{ab}
  \left(K_A^a-K_{0A}^a\right)
  \left(K_B^b-K_{0B}^b\right).
\]
Here \(K_A\in\R^3\) is the axial-vector form of the skew matrix \(A_A\):
\[
  (A_A)^a{}_b=-\epsilon^a{}_{bc}K_A^c.
\]
The tensors \(C\) and \(D\) encode stretching/shearing and bending/twisting
stiffnesses. The preferred fields \(\Gamma_0,K_0\) represent intrinsic strain
and intrinsic curvature. This formula is schematic but important: it shows how
ordinary elastic energy becomes an energy on translational strain and gauge
connection data.

\begin{remark}[Rigid body as the zero-gradient special case]
If \(R(X,t)=R(t)\) and
\[
  \varphi(X,t)=r(t)+R(t)X,
\]
then
\[
  A_A=R^{-1}\partial_A R=0,
  \qquad
  \Gamma_A=R^{-1}\partial_A\varphi=E_A.
\]
Thus a rigid body is the special Cosserat field with no spatial frame
connection and no translational strain beyond the reference directors. A
deformable Cosserat body is obtained by allowing \(R\) and \(\Gamma_A\) to vary
with \(X\).
\end{remark}

\section{Curvature, Torsion, And Defects}
\label{sec:curvature-torsion-defects}

If a connection is globally of the form
\[
  A_A=R^{-1}\partial_A R,
\]
then it is pure gauge and satisfies the Maurer-Cartan identity:
\[
  F_{AB}
  =
  \partial_AA_B-\partial_BA_A+[A_A,A_B]
  =
  0.
\]
Nonzero curvature means that local frames cannot be integrated into a single
smooth orientation field without singularities or defects. In defect theory,
curvature is associated with disclinations, or rotational defects.

To describe translational defects, introduce a local coframe \(\theta^a\), which
measures material line elements in a local frame. The connection one-forms
\(\omega^a{}_b\) are the components of the same \(\mathfrak{so}(3)\)-valued
connection written above: if the matrix entries of \(A_A\) are
\((A_A)^a{}_b\), then
\[
  \omega^a{}_b=(A_A)^a{}_b\,\dd X^A.
\]
With this notation, the torsion two-form is
\[
  T^a=\dd\theta^a+\omega^a{}_b\wedge\theta^b.
\]
Torsion measures closure failure of infinitesimal parallelograms and is
associated with dislocation density. In this language:
\[
\begin{aligned}
  \text{curvature} &\longleftrightarrow \text{rotational mismatch},\\
  \text{torsion} &\longleftrightarrow \text{translational mismatch},\\
  \text{incompatible strain} &\longleftrightarrow \text{nonintegrable local data}.
\end{aligned}
\]
The compatibility identities are the differential bookkeeping rules for these
defect densities. With covariant exterior derivative \(D\),
\[
  D F=0,
  \qquad
  D T^a=F^a{}_b\wedge\theta^b.
\]
The first identity says that disclination density satisfies a Bianchi identity.
The second says that dislocation density is not independent of rotational
curvature when the local frame itself has curvature. These identities are the
continuum-defect analogue of \(\nabla\cdot(\nabla\times A)=0\): they tell which
defect fields can be prescribed consistently.

The gauge theory is not decorative terminology. It records which descriptions
depend on arbitrary local choices of frame and which quantities are invariant.

\section{Cosserat Rod As A One-Dimensional Gauge Theory}
\label{sec:cosserat-rod}

A planar rod has material coordinate \(s\), centerline \(r(s)\in\R^2\), and
orientation angle \(\theta(s)\). Its director is
\[
  d_1=(\cos\theta,\sin\theta).
\]
For an inextensible unshearable rod,
\[
  r'(s)=d_1(s).
\]
The connection is the scalar curvature
\[
  \kappa(s)=\theta'(s).
\]
Given \(\kappa\),
\[
  \theta(s)=\theta(0)+\int_0^s\kappa(\sigma)\,\dd\sigma,
\]
and
\[
  r(s)=r(0)+\int_0^s
  (\cos\theta(\sigma),\sin\theta(\sigma))\,\dd\sigma.
\]
This is the reconstruction problem: a gauge field \(\kappa\,\dd s\) determines
shape after choosing initial position and frame.

The bending energy
\[
  E_{\mathrm{bend}}
  =
  \frac12\int_0^L B(\kappa-\kappa_0)^2\,\dd s
\]
compares actual curvature with preferred curvature. Here \(B\) is bending
stiffness and \(\kappa_0\) is intrinsic curvature. In gauge language, the
energy penalizes mismatch between a connection and a preferred connection.

\section{What To Take Forward}
\label{sec:deforming-body-take-forward}

The progression from rigid body to the gauge theory of deforming bodies is:
\[
\begin{aligned}
  \text{rigid body:}&\quad R(t),\quad R^{-1}\dot R,\\
  \text{elastic body:}&\quad \varphi(X,t),\quad F=\partial_X\varphi,\quad W(F),\\
  \text{deformable body:}&\quad g(t),\ s(t),\quad g^{-1}\dot g=-\mathcal A(s)\dot s,\\
  \text{Cosserat continuum:}&\quad R(X,t),\quad R^{-1}\partial_A R,\\
  \text{defect theory:}&\quad A,\quad F=\dd A+A\wedge A,\quad T=\dd\theta+\omega\wedge\theta.
\end{aligned}
\]
The first line is ordinary rigid-body mechanics. The second is
classical elasticity: deformation gradients, stored energy, stress, and elastic
waves. The third is Shapere-Wilczek gauge kinematics over shape space. The
fourth attaches a local rigid body to every material point. The fifth allows
the local frame and
translation data to be incompatible, producing curvature and torsion as physical
defect densities.


\section{Chapter-End Laboratories And Exercises}
\label{sec:deforming-body-chapter-end-labs}

These laboratories make the gauge language operational: a small shape loop
measures curvature, and a finite-dimensional rotor model shows how
noncommuting connection components produce body rotation.

\subsection{Mechanical Connection: Curvature From A Small Shape Loop}
\label{sec:detailed-connection-curvature}

Let \(s^1,s^2\) be two shape coordinates and let the zero-momentum
reconstruction equation be
\[
  g^{-1}\dot g
  =
  -\mathcal A_\alpha(s)\dot s^\alpha.
\]
For a tiny rectangular loop based at \(s_0\), first move by
\(\Delta s^1\), then \(\Delta s^2\), then back by \(-\Delta s^1\), then back
by \(-\Delta s^2\). To second order in the side lengths, the net group element
is
\[
  g_{\mathrm{loop}}
  =
  \exp(-\mathcal A_2\Delta s^2)
  \exp(-\mathcal A_1\Delta s^1)
  \exp(\mathcal A_2\Delta s^2)
  \exp(\mathcal A_1\Delta s^1)
\]
plus the correction from the variation of \(\mathcal A_\alpha(s)\) across the
loop. The Baker-Campbell-Hausdorff formula gives the commutator contribution,
while the spatial variation gives derivative terms. The result is
\[
  g_{\mathrm{loop}}
  =
  \exp\left(
    -\mathcal F_{12}\Delta s^1\Delta s^2
    +
    O(|\Delta s|^3)
  \right),
\]
where
\[
  \boxed{
  \mathcal F_{12}
  =
  \partial_1\mathcal A_2-\partial_2\mathcal A_1
  +
  [\mathcal A_1,\mathcal A_2].
  }
\]
This is why curvature is the local density of geometric phase. Even if the
connection components are constant, a nonabelian structure group can give
nonzero holonomy through the commutator term. The helicopter rotor laboratory
later in this chapter is exactly of this type.

\subsection{Helicopter Rotor Model As A Gauge Connection}
\label{sec:lab-helicopter-gauge}

A toy helicopter model makes the Shapere-Wilczek connection finite-dimensional.
It is the most compact laboratory version of the inertia-operator connection
derived in Section~\ref{subsec:explicit-gauge-equation}: the general shape
coordinate \(s\) is replaced by two rotor angles \((\alpha,\beta)\), and the
zero-momentum equation can be solved by inspection. The same calculation is
presented in its stripped-down conceptual form in
Section~\ref{sec:two-rotor-deforming-body-model}; the laboratory here is the
worked-problem version, with the kinetic energy and boundary of the shape
cycle written out explicitly.

The calculation has one main checkpoint. The rotor angles return to their
initial values after a closed stroke, so any final rotation of the body cannot
be attributed to a changed shape. It must come from the noncommuting
connection components. This is the finite-dimensional version of geometric
phase.

Let the body have isotropic base inertia
\[
  I_{\mathrm{body}}=I_0\mathbf 1.
\]
Let one internal rotor spin about the body \(z\)-axis with angle \(\alpha\) and
inertia \(I_1\), and let a second rotor spin about the body \(y\)-axis with
angle \(\beta\) and inertia \(I_2\). The internal shape space is the torus
\[
  \mathcal S=S^1_\alpha\times S^1_\beta.
\]
Let \(\Omega\in\R^3\) be the body angular velocity. In the body frame, the
rotor angular velocities are
\[
  \Omega+\dot\alpha e_z,
  \qquad
  \Omega+\dot\beta e_y.
\]
Keeping only the rotational kinetic energy relevant for reconstruction,
\[
  T=
  \frac12I_0\norm{\Omega}^2
  +
  \frac12I_1(\Omega_z+\dot\alpha)^2
  +
  \frac12I_2(\Omega_y+\dot\beta)^2.
\]
The body angular momentum conjugate to \(\Omega\) is
\[
  \Pi
  =
  I_0\Omega
  +
  I_1(\Omega_z+\dot\alpha)e_z
  +
  I_2(\Omega_y+\dot\beta)e_y.
\]
Set total angular momentum to zero. Componentwise,
\[
  (I_0+I_1)\Omega_z+I_1\dot\alpha=0,
\]
\[
  (I_0+I_2)\Omega_y+I_2\dot\beta=0,
\]
\[
  I_0\Omega_x=0.
\]
Therefore
\[
  \Omega
  =
  -\frac{I_1}{I_0+I_1}\dot\alpha\,e_z
  -
  \frac{I_2}{I_0+I_2}\dot\beta\,e_y.
\]
In connection notation,
\[
  g^{-1}\dot g
  =
  -\mathcal A_\alpha\dot\alpha-\mathcal A_\beta\dot\beta,
\]
with
\[
  \mathcal A_\alpha=
  \frac{I_1}{I_0+I_1}e_z,
  \qquad
  \mathcal A_\beta=
  \frac{I_2}{I_0+I_2}e_y.
\]
The components are constant, but the structure group is nonabelian:
\[
  [e_z,e_y]\neq0.
\]
Thus the curvature contains the commutator term
\[
  \mathcal F_{\alpha\beta}
  =
  [\mathcal A_\alpha,\mathcal A_\beta].
\]
A small rectangular cycle in \((\alpha,\beta)\) can therefore produce a net
body rotation even though each rotor angle returns to its original value. This
is the finite-dimensional version of falling-cat holonomy.

\begin{figure}[htbp]
\centering
\begin{tikzpicture}[scale=1.0, >=Latex, font=\small]
  \begin{scope}
    \draw[->] (-0.2,0) -- (4.4,0) node[right] {\(\alpha\)};
    \draw[->] (0,-0.2) -- (0,3.1) node[above] {\(\beta\)};
    \fill[blue!8] (0.85,0.65) rectangle (3.25,2.35);
    \draw[blue!70!black, very thick, -{Latex[length=3mm]}] (0.85,0.65) -- (3.25,0.65)
      node[midway, below] {\(\Delta\alpha\)};
    \draw[blue!70!black, very thick, -{Latex[length=3mm]}] (3.25,0.65) -- (3.25,2.35)
      node[midway, right] {\(\Delta\beta\)};
    \draw[blue!70!black, very thick, -{Latex[length=3mm]}] (3.25,2.35) -- (0.85,2.35)
      node[midway, above] {\(-\Delta\alpha\)};
    \draw[blue!70!black, very thick, -{Latex[length=3mm]}] (0.85,2.35) -- (0.85,0.65)
      node[midway, left] {\(-\Delta\beta\)};
    \node[align=center] at (2.05,3.35) {closed shape loop\\on the rotor torus};
  \end{scope}
  \begin{scope}[xshift=5.4cm]
    \draw[->, thick] (0,0) -- (1.6,0) node[right] {\(e_x\)};
    \draw[->, thick] (0,0) -- (0,1.6) node[above] {\(e_z\)};
    \draw[->, thick] (0,0) -- (-1.1,-0.8) node[left] {\(e_y\)};
    \draw[red!75!black, thick, ->] (0.95,0.48)
      arc[start angle=90, end angle=-230, x radius=0.22, y radius=0.48];
    \node[red!75!black, fill=white, inner sep=1pt] at (1.55,0.55)
      {rotation about \(e_x\)};
    \node[black!65, fill=white, inner sep=1pt] at (0.72,1.95)
      {body-frame axes};
    \node[red!75!black, align=center, fill=white, inner sep=1.5pt] at (0.92,-1.55)
      {net body rotation\\from \([\mathcal A_\alpha,\mathcal A_\beta]\)};
  \end{scope}
\end{tikzpicture}
\caption{Helicopter-rotor gauge laboratory. The connection components are
constant, but they point in noncommuting body directions. A closed rectangle in
shape space therefore has curvature from the commutator term and can produce a
net body rotation.}
\label{fig:lab-helicopter-gauge-loop}
\end{figure}
\FloatBarrier

\section{Associated Demos}
\label{sec:deforming-body-demos}

\begin{itemize}
\item \path{demos/python/cosserat_rod_demo.py}, a planar frame reconstruction
from curvature.
\item \path{demos/python/deforming_body_gauge.py}: finite-dimensional
two-rotor gauge holonomy.
\item \path{demos/mathematica/GaugeDeformingBody.wl}, a finite-dimensional
gauge-connection toy model.
\item Suggested extension: a low-Reynolds swimmer whose connection is computed
from a Stokes boundary-value problem.
\end{itemize}
